%% 04_contribution.tex — 논문 작성의 기여도
%% ✏️ 아래에 각 저자의 기여 내용을 입력해 주세요.
%%
%% 사용 방법:
%%   \contribA{연번}{논문 위치}{기여 내용}   ← 저자 A의 기여 행 추가
%%   \contribB{연번}{논문 위치}{기여 내용}   ← 저자 B의 기여 행 추가
%%
%% 예시:
%%   \contribA{1}{2.1. 서론}{연구 배경을 조사하고 서론을 작성하였다.}
%%   \contribA{2}{3. 본론}{핵심 정리의 증명을 담당하였다.}
%%   \contribB{1}{2.2. 이론적 배경}{관련 선행 연구를 정리하였다.}
%%
%% '논문 위치'에는 해당 기여가 반영된 절 번호(예: 3.2.)를 씁니다.
%% 행 개수는 저자마다 다르게 해도 됩니다.
%% ─────────────────────────────────────────────

% ── 저자 A ────────────────────────────────────
%TODO: 저자 A의 기여 내용을 아래에 입력하시오.
\contribA{1}{}{}


% ── 저자 B ────────────────────────────────────
%TODO: 저자 B의 기여 내용을 아래에 입력하시오.
\contribB{1}{}{}
